% =====================================================================
%  CARNET D'ENTRAÎNEMENT — FICHE 8 — THÈME 1 — NIVEAU DIFFICILE — ÉNONCÉS
% =====================================================================
\documentclass[11pt,a4paper]{article}
\usepackage[utf8]{inputenc}
\usepackage[T1]{fontenc}
\usepackage{lmodern}
\usepackage[french]{babel}
\usepackage{amsmath,amssymb,mathtools}
\usepackage{icomma}
\usepackage{siunitx}
\sisetup{output-decimal-marker={,},per-mode=power,group-separator={\,},group-minimum-digits=5}
\usepackage[version=4]{mhchem}
\usepackage{xcolor}
\usepackage{tikz}
\usepackage[most]{tcolorbox}
\usepackage{enumitem}
\usepackage{array}
\usepackage{fancyhdr}
\usepackage[hidelinks]{hyperref}
\usetikzlibrary{arrows.meta,positioning,calc}
\usepackage[a4paper,margin=2.1cm,headheight=26pt,headsep=10pt]{geometry}
\usepackage{titlesec}

\definecolor{primary}{RGB}{150,98,162}
\definecolor{primarydark}{RGB}{92,52,110}
\definecolor{exocol}{RGB}{200,60,40}
\definecolor{travail}{RGB}{30,110,180}
\definecolor{chaleur}{RGB}{214,72,33}

\tcbset{boxbase/.style={enhanced,breakable,boxrule=0.9pt,arc=2.5pt,left=3mm,right=3mm,top=2.3mm,bottom=2.3mm,fonttitle=\bfseries,coltitle=white,toptitle=0.6mm,bottomtitle=0.6mm}}
\newtcolorbox[auto counter]{exercice}[1][]{boxbase,colback=primary!4,colframe=primary,title={\textsc{Exercice}~\thetcbcounter\ \ifx\relax#1\relax\relax\else\textmd{--- \itshape #1}\fi}}

\newcommand{\dH}{\Delta H}
\newcommand{\dU}{\Delta U}

\pagestyle{fancy}
\fancyhf{}
\fancyhead[L]{\small\textcolor{primarydark}{\textbf{Fiche 8} — Thème 1 : Premier principe et enthalpie}}
\fancyhead[R]{\small\textcolor{primarydark}{\textsc{Difficile}}}
\fancyfoot[C]{\small\thepage}
\renewcommand{\headrulewidth}{0.8pt}
\renewcommand{\headrule}{\hbox to\headwidth{\color{primary}\leaders\hrule height \headrulewidth\hfill}}
\setlength{\parindent}{0pt}\setlength{\parskip}{3pt}

\begin{document}

\begin{center}
\begin{tikzpicture}
\node[rectangle,rounded corners=7pt,fill=primary,text=white,minimum width=16cm,
      minimum height=1.1cm,align=center,inner sep=6pt]
  {{\large\bfseries Thème 1 — Premier principe et enthalpie}\\[1pt]{\small Niveau \textsc{Difficile} — énoncés}};
\end{tikzpicture}
\end{center}

\begin{exercice}[détente d'un gaz dans un piston — bilan complet]
Un cylindre fermé par un piston mobile contient un gaz maintenu à la \textbf{pression extérieure
constante} $p = \SI{1.0e5}{\pascal}$. On chauffe le gaz : il \textbf{reçoit} $Q = \SI{+250}{\joule}$ de
chaleur et se dilate de $V_{\text{i}} = \SI{1.0}{\litre}$ à $V_{\text{f}} = \SI{3.0}{\litre}$.
\begin{center}
\begin{tikzpicture}[scale=0.9]
  \draw[thick] (0,0) rectangle (3,1.6);
  \fill[primary!12] (0,0) rectangle (1.3,1.6);
  \draw[thick,fill=gray!35] (1.3,0) rectangle (1.55,1.6);
  \draw[->,travail,line width=1pt] (1.7,0.8) -- (2.7,0.8) node[right,font=\scriptsize,travail]{piston};
  \draw[->,chaleur,line width=1pt] (-0.9,0.8) -- (-0.1,0.8) node[midway,above,font=\scriptsize,chaleur]{$Q$};
  \node[font=\scriptsize] at (0.65,0.8){gaz};
\end{tikzpicture}
\end{center}
\begin{enumerate}[label=\alph*),leftmargin=2em,topsep=2pt,itemsep=2pt]
  \item Calculer le travail $W$ échangé par le gaz (avec son signe). Le gaz reçoit-il ou cède-t-il ce travail ?
  \item En déduire la variation d'énergie interne $\dU$.
  \item L'énergie interne a-t-elle augmenté ou diminué ? Interpréter : que devient la chaleur reçue ?
\end{enumerate}
\end{exercice}

\begin{exercice}[combustion du carbone — énergie et quantité de matière]
La combustion complète du carbone à pression constante s'écrit
\;$\ce{C}(s) + \ce{O2}(g) \rightarrow \ce{CO2}(g)$,\; avec $\dH = \SI{-394}{\kilo\joule\per\mol}$.
On donne $M(\ce{C}) = \SI{12}{\gram\per\mol}$.
\begin{enumerate}[label=\alph*),leftmargin=2em,topsep=2pt,itemsep=2pt]
  \item Décrire le diagramme enthalpique : lequel, des réactifs ou des produits, est le plus haut ?
        Quel est le signe de $\dH$, et la réaction est-elle exo ou endothermique ?
  \item Quelle énergie thermique est dégagée par la combustion de \SI{3.0}{\mol} de carbone ?
  \item On brûle une masse $m$ de carbone et on récupère \SI{1970}{\kilo\joule}. Calculer la quantité de
        matière brûlée, puis la masse $m$.
  \item La combustion \emph{incomplète} \;$\ce{C}(s) + \tfrac12\,\ce{O2}(g) \rightarrow \ce{CO}(g)$\; a
        $\dH = \SI{-111}{\kilo\joule\per\mol}$. Pour une même mole de carbone, laquelle des deux
        combustions dégage le plus d'énergie ? Commenter.
\end{enumerate}
\end{exercice}

\begin{exercice}[distinguer $\dU$ et $\dH$ — l'ordre de grandeur du travail]
Une réaction en phase gazeuse à \textbf{pression constante} $p=\SI{1.0e5}{\pascal}$ s'accompagne d'un
\textbf{dégagement de gaz} : le volume augmente de $\Delta V = \SI{+2.0}{\litre}$. La chaleur mesurée à
pression constante vaut $Q_p = \dH = \SI{-180}{\kilo\joule}$.
\begin{enumerate}[label=\alph*),leftmargin=2em,topsep=2pt,itemsep=2pt]
  \item Calculer le travail $W$ (en \si{\joule}, puis en \si{\kilo\joule}).
  \item En déduire $\dU$ à l'aide de $\dU = \dH + W$.
  \item Comparer les valeurs de $\dU$ et de $\dH$. Le travail des forces de pression est-il important
        devant la chaleur de réaction ? Que peut-on retenir de général ?
\end{enumerate}
\end{exercice}

\end{document}
